\section{Loop categories over moving divisors}\label{sec:loop}

The local comparison will involve marked points that move and collide.  We
therefore construct its categories and pairings over a parameter scheme,
including their infinitesimal transport.  The essential output is the
pairing \eqref{eq:L3} and its continuous, parameter-linear right adjoints.
The construction also identifies the strong equivariance used in that
pairing with the Harish--Chandra models used for the representation-theoretic
calculations in the next sections.

\subsection{Divisors, determinant lines, and parameter directions}

Let $S$ be a smooth affine scheme with a finite tuple of maps $x_l:S\to X$.
Additional parameters are allowed.  Complete along the \emph{sum} of their
graphs.  Repeating a graph changes its multiplicity but not the resulting
formal completion or meromorphic algebra.  After an \'etale localization,
separate disjoint clusters and choose a curve coordinate $z$ and a frame
of the theta characteristic.  On one cluster the divisor has equation
\[
                 u=P(z)=\prod_l(z-x_l).
\]
Only the formal neighborhood of the marked points in this coordinate
chart is used.  Write $LG$ and $H=L^+G$ for the resulting relative loop and
arc groups.  The groups are twisted by $P_0$; the displayed coordinates
include a trivialization of $P_0$.  In particular, the intrinsic unipotent
group is $N_0$.

The relative determinant extension is the multiplicative line
\[
     \mathcal L(g)=\det(g\mathfrak g_{\mathcal O}:\mathfrak g_{\mathcal O}),
\]
with first lattice over second lattice and the adjoint action on
$\mathfrak g$.  The lattice indices vanish, so this is an ordinary
multiplicative line.  It is the local line in the global Hecke
correspondence.  Its Lie cocycle is
\[
                       (x,y)\longmapsto
                       \operatorname{Res}\operatorname{Kil}(dx,y).
\]
Indeed, the Tate trace cocycle is the trace of the two off-diagonal
blocks relative to the regular lattice; on modes $az^m,bz^{-m}$, with
$m>0$, it is $m\operatorname{Tr}(ab)$.  This proves the formula at distinct
points.  Both sides are polynomial identities in each bounded collection
of Laurent coefficients, and the relative residue formula gives the same
identity on the collision locus.

Put $\ell=\widehat{\mathfrak g}_{\mathcal K}$ for this Tate Lie algebra
with its central extension.  Its central element acts by
\[
                         -s=\frac{k_a}{2h^\vee}.
\]
Thus the basic-form level is $k_a$ and the shifted level is $a$.  The
sign comes from the frame convention: a frame for a modification glues
from the inside to the outside, and sections transform by its inverse
power.  Left convolution consequently uses the twist $\mathcal L^{-s}$.

For a module $M$, a loop $g$, and a determinant lift $\widetilde g$, the
family of conjugated modules, paired with this convolution twist, is
given in left-module normalization by
\begin{equation}\label{eq:L1}
 \begin{split}
 x&\longmapsto\operatorname{act}_M(\operatorname{Ad}_{g^{-1}}x),\\
 \partial_v&\longmapsto\partial_v+
       \operatorname{act}_M(\widetilde g^{-1}\partial_v\widetilde g).
 \end{split}
\end{equation}
Here $\partial_v$ differentiates in the parameter direction while holding
the curve coordinate $z$ fixed.  Changing the determinant lift adds
$-s$ times its logarithmic derivative, so this family has twist $s$.
The second line is essential: conjugation alone does not specify the
connection when the marked points move.  We construct the categorical
meaning of \eqref{eq:L1} below.

We use crystals of presentable categories in the sense of
\cite[Section~2 and Appendix~B, especially Sections~B.3.8--B.3.11
and~B.14]{GLCII}.  By
\cite[Theorem~2.6.3]{Gaitsgory1Affine}, a crystal over $S$ is equivalently
a $D(S)$-module category.  Its underlying quasi-coherent category is
\[
 \mathcal E^0=\operatorname{QCoh}(S)
                   \otimes_{D(S)}\mathcal E(S),
 \qquad b:\mathcal E(S)\longrightarrow\mathcal E^0.
\]
The functor in this tensor product is the left-crystal forgetful functor:
it sends $F^{\langle S\rangle}$ to its underlying quasi-coherent module.
After applying $b$, equivariant operations on smooth schemes over $S$
are relative operations.  This follows either from crystal base change
or from the relative differential-operator algebra.

We will check crystal compatibilities on the formal diagonal groupoid
of $S$.  This means all finite infinitesimal neighborhoods of every
multiple diagonal, with their coherent pullback maps.  Since $S$ is
smooth, the projections may be computed using cofinal neighborhoods
finite and flat over $S$.  The use of this groupoid will permit moving
congruence subgroups; it does not declare a fixed finite jet group to
be horizontally constant.

\subsection{Relative convolution and equivariance}

Let
\[
                     \mathcal A=D_{*,\mathcal L^{-s}}(LG).
\]
The subscript specifies star convolution and its pro convention.
In the ind direction we use closed direct images.  In the pro direction
we use systems compatible under star direct image.  The pro projections
have smooth unipotent affine-space fibers; their star-compatible limit
is also the colimit under normalized smooth pulls, the left adjoints of
those direct images.

Concretely, on a bounded part of $LG$ stable on the right under $H$, take
quotients by sufficiently deep congruence subgroups $K\subset H$.  Left
quotients give the same category: bounded conjugation sends sufficiently
deep congruences into any specified congruence, so both presentations
have common higher-jet refinements.  A bounded part may be the preimage
of a closed bound in the affine Grassmannian.  Multiplication is star
direct image with the multiplicative line.  All products are relative:
for the finite-type schemes used in these presentations,
\[
                 D(U)\otimes_{D(S)}D(V)
                       \simeq D(U\times_S V).
\]
This is the D-module tensor-product formula, followed by base change
along the base diagonal; it also holds with the indicated twists.
For any finite string of products and any output jet level, sufficiently
deep input levels make the calculation finite-dimensional.  Composition
of star direct images then proves associativity, including its higher
coherences.  The delta system at the identity is the unit.

We need an equivariance calculation over $S$, where jet groups can vary
with the divisor.  A \emph{strong action} in the following statement
means an action of the convolution category of D-modules.  This is
stronger than an algebraic action on the underlying category.

\begin{proposition}[Relative invariant and coinvariant calculation]
\label{loop:relative-equivariance}
Let $J\to S$ be one of the smooth affine finite jet groups of the arc,
congruence, or unipotent groups used here.  Put
$A_J=D_*(J)$ and let $\mathbf1_J=D(S)$ be its module given by relative
de Rham direct image.  For every $D(S)$-linear strong $J$-category
$\mathcal E$, the norm calculation identifies
$\mathcal E_J$ with $\mathcal E^J$, naturally also in additional
coefficient categories.  Under this duality:
\begin{enumerate}
\item geometric relative tensor products over $A_J$ are D-modules on
contracted products, with star integration in a free group variable;
\item $\operatorname{End}_{A_J}(\mathbf1_J)=D(B_SJ)$, with ! tensor;
\item the unbalanced evaluation
\[
 \operatorname{ev}_J:
 {}_{A_J}\mathbf1_J\otimes_{D(S)}\mathbf1_J{}_{A_J}
       \longrightarrow A_J
\]
has a continuous right adjoint, linear for the left and right
$A_J$-actions.  Both actions on $\mathbf1_J$ are given by the
de Rham augmentation.
\end{enumerate}
The assertions hold for the pro-smooth groups in use, whose kernels
over a finite jet stage are pro-unipotent, and commute with change of
smooth parameter base.  Character versions use the corresponding
exponential and its inverse on opposite sides.
\end{proposition}

\begin{proof}
For a constant finite-dimensional group, this is strong-group-action
duality; see the distinction between group schemes and loop indgroups
in \cite[Section~1.5.4]{RaskinW}.  We reduce the relative case to this
one without assuming that $J$ is constant.

Locally on $S$, embed $J$ in a constant general linear group $D_0$ over
$S$, with schematic quotient $W=D_0/J$.  The embeddings needed here
can be constructed as follows.  An arc jet group acts faithfully on a
finite free module over the finite divisor algebra.  Restriction of
scalars embeds it in a general linear group over $S$.  The intermediate
quotient parametrizes module structures over the divisor algebra for which it is
locally free of the specified rank over the divisor algebra.  This is
an open locus in an affine representation scheme.  The remaining
quotient is affine, since the original reductive group has affine
quotient in a general linear group and restriction of scalars along a
finite flat map preserves affines.  Split unipotent jet groups admit
upper-triangular embeddings using their root and congruence filtrations.
The same construction applies after restricting to a base chart.

The relative D-module category of a finite-type schematic map $U\to S$
is dualizable over $D(S)$: its coevaluation is the relative diagonal
kernel $\Delta_*\omega_U$, and evaluation is ! tensor followed by star
integration.  The triangle identities are base change on the relative
square.  Transposing an action using this duality gives its !-coaction
and hence the usual descent diagram for strong invariants.

Consider
\[
            \mathcal B=A_{D_0}\otimes_{A_J}\mathcal E.
\]
We first show, by a free-action computation, that
\[
       \mathcal B\simeq
          \bigl(D(D_0)\otimes_{D(S)}\mathcal E\bigr)^J.
\]
Here $J$ acts by right translation on $D_0$, and inversion identifies
the required right and left conventions.  If $m\mapsto m^h$ denotes
the transposed coaction, the comparison is star averaging: on a
trivialized right torsor $U\to W$, its ! family at a section $x$ is
the star integral in $h$ of
$K(xh)\otimes^!m^h$.  The coaction identity makes this an equivariant
family.  Shearing two group variables shows that convolution on $K$
is the action on $m$, so the comparison is balanced; the same
calculation for strings supplies the coherent balance.

For a trivial torsor, pulling the family to a section splits its
descent bar.  The extra degeneracy identifies its invariants with
coefficients on $W$, and the ! fiber of the displayed averaging
formula is exactly the original action.  Thus the comparison is an
equivalence on a trivial torsor.  A smooth torsor has sections \'etale
locally.  Restriction to such a cover commutes with the descent limits,
by the relative dualizability just proved, and with the tensor bars.
Descent for $D(W)$-module categories follows from the 1-affineness of
$W_{\mathrm{dR}}$.  It therefore gives the claimed equivalence over
$W$, and over $S$.  The construction is $A_{D_0}$-linear.

Associativity of relative tensors, constant-group duality, and this
free-action formula now give
\[
        \mathcal E_J\simeq\mathcal B_{D_0}
                    \simeq\mathcal B^{D_0}
                    \simeq\mathcal E^J.
\]
For the last equivalence, interchange the two commuting descent
diagrams and apply translation descent on $D_0$.  These computations
also work with coefficient categories and hence identify the claimed
module duality, rather than just its value on one category.

On the regular category, translation descent identifies the two
trivial modules with $D(S)$.  Evaluation is consequently the smooth
pull kernel on $J$, up to an invertible normalization from the base.
Its right adjoint is star integration with the corresponding smooth
shift.  The projection formula and the shear map identify the
canonical adjunction-linearity maps for the two convolutions with
isomorphisms.  The contracted-product formula and the description of
the endomorphism category follow from the same descent calculation.

Finally, for a pro-unipotent kernel, normalized smooth pull of the
unit is a convolution idempotent at every finite quotient.  Its image
is strong equivariance, and its continuous right adjoint is star
averaging.  The delta unit in the pro-convolution convention is the
filtered colimit of these idempotents for decreasing kernels: this
is checked by projection to every finite level.  This reduces the
assertions to a finite quotient.  In particular $D(B_SJ)$ may be
computed at any finite stage with pro-unipotent kernel.  All maps in
the argument are relative pull, star integration, or tensor bars,
so the construction commutes with the asserted base changes.
\end{proof}

This proposition is used only for finite-dimensional or compact
pro-groups.  Coinvariants for the indgroup $LN$ will instead be formed
as a colimit over compact subgroups.

\subsection{Renormalized Harish--Chandra categories}

We next describe the underlying quasi-coherent categories, after applying
$b$.  Every compact-generation assertion in this subsection is made
in these categories.  It does not assert generation by horizontal
objects; in particular, multipoint Kazhdan--Lusztig categories need
not be ULA-generated \cite[Remark~B.14.25]{GLCII}.  Put
$R=\mathcal O(S)$.  Modules for $\ell$ are smooth: every vector is
annihilated by a sufficiently deep Lie lattice.  Let $J$ be an arc
group, a congruence subgroup, or a pro-unipotent compact open of the
kind above, with Lie algebra $j$ and a specified splitting of the
extension.  A character $\chi_J$ is incorporated by the augmentation
\[
       K\longmapsto p_{J,*}
                      \bigl(K\otimes^!\operatorname{Exp}(\chi_J)\bigr).
\]
The corresponding Harish--Chandra condition identifies the differential
of the $J$-action with the $\ell$-action \emph{minus} $\chi_J$.
This sign agrees with \eqref{eq:L1}.

Define the renormalized category for $(\ell,J)$ as the Ind category
of the idempotent-complete stable category generated by
\begin{equation}\label{eq:L2}
                       \operatorname{Ind}_{j}^{\ell}(V),
\end{equation}
where $V$ ranges over algebraic $J$-representations that are vector
bundles over $S$.  Impose the specified central value and character
correction.  The mapping complexes on these generators are the
ordinary derived Harish--Chandra mapping complexes.  By induction
adjunction, they are the derived algebraic $J$-maps from $V$ to the
restricted target.  Renormalized pro-group representations are defined
in the same way, omitting Lie induction.

\begin{proposition}\label{loop:hc-renormalization}
The underlying quasi-coherent categories have the following properties.
\begin{enumerate}
\item The comparison with the ordinary derived category of smooth
Harish--Chandra modules is t-exact and is an equivalence on complexes
bounded below in cohomological degrees.
\item The category for $(\ell,J)$ is monadic over renormalized
$J$-representations, with the PBW induction monad.  Its t-structure is
created by the forgetful functor.
\item The compact-generator mapping calculations commute with smooth
affine base change, restriction to smooth closed collision and graph
charts, residue fields, and the finite flat scalar
extensions used on infinitesimal diagonals.  In the last case the
generators are perfect relative to the scalar extension.
\item Restriction to a smaller compact open with smooth finite-dimensional
quotient preserves compact objects.  Between pro-unipotent opens it
is fully faithful and has a continuous right adjoint.
\end{enumerate}
\end{proposition}

\begin{proof}
The pointwise renormalization and bounded comparison are developed in
\cite[Sections~9.11--9.13 and~10.8]{RaskinHomological}; we give the
relative calculation needed here.  First consider finite group stages.
The quotient presentations in
the proof of Proposition~\ref{loop:relative-equivariance} give the
resolution property.  Indeed, quasi-coherent representations are
generated by vector bundles: use generation on the quasi-affine
homogeneous space, and then finite algebraic subrepresentations of
its sections.  General linear group invariants are exact in
characteristic zero, and a quasi-affine scheme of finite type has
bounded cohomological dimension.  On the smooth base, bounded coherent
equivariant complexes are therefore perfect and compact.  Ordinary
equivariant descent supplies the derived representation category at
these stages.  Restricted perfects generate after the stated base
changes, either by the same quotient presentation or by testing after
the affine pushforward.

Every algebraic representation of a pro-group is exhausted by the
subobjects fixed by deeper kernels, as follows from its coaction.
Mappings between bounded coherent objects are computed by the group
cobar resolution.  If the target complexes are bounded below in a
fixed range, maps from a finite vector bundle commute with filtered
colimits: degree by degree the cobar computation uses only a finite
part of the resulting first-quadrant total complex.  This proves the
bounded-below comparison for renormalized group representations.
More explicitly, the hearts are generated by the finite-stage
representations; bounded complexes lift using the mapping calculation;
and bounded-below complexes are obtained as colimits of their good
upper truncations.  The same calculation shows that the renormalized
pro-category is the colimit of the finite-stage renormalized
categories.  The finite-stage transition functors need not be fully
faithful.

For later use, in the pro-unipotent case the compact representation
category is thickly generated over $R$ by the trivial representation.
At a finite stage, the successive vector-group quotients act by
commuting locally nilpotent operators.  On a coherent representation
they supply a finite filtration by coherent subobjects with trivial
successive action; repeat with the quotient group and resolve over
the smooth affine base.  This proves the assertion at each stage and
hence in the pro-category.

The Lie induction monad is the PBW monad.  The quotient of the Tate
Lie algebra by an open lattice is a union of finite locally free
summands, so PBW makes induction exact.  The adjoint $J$-action is
algebraic: sufficiently high congruences fix each specified bounded
Laurent expression and each vector fixed at a sufficiently deep
level.  The monad thus acts on the representation categories just
constructed.  Its extension to their Ind categories agrees on
bounded-below objects with the ordinary induction monad.  Its free
modules have exactly the mapping complexes prescribed in
\eqref{eq:L2}, proving the monadic assertion.  The ordinary induction
and forgetful functors are exact, the latter is conservative and
commutes with colimits, and all monad iterates agree in the
bounded-below comparison.  This proves the first two assertions for
Harish--Chandra modules.

For base change, compute the generator mappings with the same relative
cobar complexes and PBW modules.  Their terms are $R$-flat, and the
asserted changes have finite Tor dimension; for a regular closed
chart this is given by its finite Koszul resolution.
The bounded-below calculation therefore commutes with the derived
tensor product.  A finite flat algebra from an infinitesimal diagonal
is treated as an algebra object in these categories.  Its compact
generators are obtained from perfect scalar extensions, not by
declaring every coherent module over that possibly singular algebra
compact.  The free-generator cobar calculation gives the same
base-change and bounded-below assertions there.

Finally, restriction of \eqref{eq:L2} to a smaller open is resolved
by the finite-dimensional relative Chevalley complex for the
difference of the Lie lattices.  Retain the smaller group's
equivariance throughout.  This finite resolution proves compact
preservation.  The opens may first be refined through a common
deeper congruence with smooth quotients.  If both opens are
pro-unipotent and the smaller is normal, the assertion of full
faithfulness is the equality of the relative Chevalley calculation
and algebraic group cohomology for the finite-dimensional unipotent
quotient.  For a nonnormal inclusion use a still smaller normal
subgroup.

On bounded-below modules the essential image consists of those
objects whose corrected Lie action integrates to the larger
pro-unipotent group.  Such integration is unique and is preserved by
extensions.  Smoothness reduces this assertion, vector by vector,
to local nilpotence for a finite-dimensional unipotent quotient.
The fully faithful assertion then extends from the compact
presentations to the Ind categories.  Compact preservation implies
that the right adjoint to restriction is continuous.
\end{proof}

We define the renormalized Kac--Moody category $\mathcal M^0$ as the
colimit of these categories over decreasing untwisted pro-unipotent
congruence levels.  The maps into this category will be called
forgetful maps.  They are distinct from passage to the ordinary
abelian module model.  In particular, Proposition~\ref{loop:hc-renormalization}
does not identify the two unbounded categories; this restriction is
essential for Whittaker objects, as explained in
\cite[Sections~1.5.7--1.5.9]{RaskinW}.

\subsection{Infinitesimal transport and the strong loop action}

The preceding categories must now be supplied with their connection
in $S$.  The issue is that a fixed congruence level moves when a graph
moves.  The full completed arc algebra does not have this defect.

\begin{proposition}\label{loop:crystal}
The categories just constructed and their restriction and induction
functors descend along the formal diagonal groupoid of $S$.
Their sections are $D(S)$-linear categories, with underlying categories
as above.  The construction commutes with smooth parameter change and
with separation into disjoint clusters, including after restriction
to a collision locus.
\end{proposition}

\begin{proof}
Two infinitesimally close lifts of the marked graphs give cofinal
divisor ideals.  Explicitly, if their monic equations $P,Q$ agree
modulo an ideal $I$ with $I^r=0$, the binomial expansion gives
\[
               Q^{n+r-1}\in(P^n),\qquad
               P^{n+r-1}\in(Q^n).
\]
Every surviving term in the first expansion contains at least $n$
factors of $P$, and the other inclusion is symmetric.  They therefore
define the same completions,
full arc groups, and Laurent algebras.  The determinant extension
has the same canonical identification.

Choose a congruence deep enough for one lift to lie in the two
levels being compared.  It may be chosen normal in the full arc
group.  The two relative quotients are smooth: the quotient polynomial
in each such divisibility is monic, so the finite difference of the
corresponding lattices is locally free even over the nonreduced base.
Restriction to this common deeper level is fully faithful and
compact-preserving by Proposition~\ref{loop:hc-renormalization}, also
after the finite flat scalar extensions under consideration.

The two essential images there agree.  To see this on bounded-below
modules, filter by the powers of the nilpotent base ideal.  On every
successive quotient the two integrability conditions coincide.
Integrability is stable under extensions, so it coincides on the
whole module.  Equivalently, a change with coefficients in a
nilpotent ideal integrates by the finite exponential.  Restriction
is t-exact, by the group-representation comparison followed by the
induction monad.  Thus the same argument applies to the scalar
extensions of the compact generators.  Their images are compact in
the common-depth category; when they lie in the other fully faithful
subcategory they remain compact there.  The two compact subcategories
therefore agree, and so do their Ind-completions.

These identifications are independent of the common depth, because
passing to a deeper one is fully faithful.  On multiple
infinitesimal diagonals choose a single common deeper level for
all lifts.  All identifications are then restrictions of the same
full Laurent algebra, and their compositions agree there.  This
supplies the higher cocycle compatibilities, not just transport on
first-order tangent vectors.  Smooth descent and the formal
diagonal description of crystals give the claim.  The base-change
mapping calculations in Proposition~\ref{loop:hc-renormalization}
identify the underlying categories on a new smooth parameter
scheme with its corresponding Harish--Chandra categories.

The derivative at fixed $z$ preserves the full regular arc
algebra.  It need not preserve a chosen finite-depth subgroup;
the comparison through deeper levels is exactly what makes that
derivative act on the crystal.  Disjoint-cluster products follow
from products of the completed algebras.  The cofinal-ideal
argument applies equally at collisions.
\end{proof}

Write $\mathcal M$ for the resulting crystal category of Kac--Moody
modules.  We next verify that its geometric loop action is the strong
action whose invariants were used above.

\begin{proposition}\label{loop:strong-action}
Formula~\eqref{eq:L1}, paired by ! tensor and star integration, defines
a unital associative $\mathcal A$-action on $\mathcal M$ in
$D(S)$-module categories.  For each compact open $J$ used above, its
strong $(J,\chi_J)$-equivariant category is the renormalized
Harish--Chandra category for $(\ell,J)$ with differential
$\operatorname{act}_\ell-\chi_J$.  These identifications respect
the forgetful maps into $\mathcal M$.
\end{proposition}

\begin{proof}
Suppose an input belongs to a congruence level $K$, and take a bounded
right-stable part $Z$ of the loop group.  Project a convolution kernel
to $Z/K$.  On this quotient it is paired with $M^g$ from
\eqref{eq:L1}.  A sufficiently deep output level $K'$ satisfies
$g^{-1}K'g\subset K$ uniformly on the bound; thus $M^g$ has values in
the smaller Harish--Chandra category, and the pairing and integration
are ordinary relative operations with that category as coefficients.

One can compute $M^g$ after ! pull to any smooth affine plot of $Z/K$.
Choose a local loop lift and determinant frame.  At $b$, the first
line of \eqref{eq:L1} is conjugation restricted to $K'$.  The
positive-depth splittings agree with conjugation after increasing
the depth; this follows directly from the relative determinant
calculation.  Across an infinitesimal diagonal, two lifts differ by
a formal group element.  Exponentiating its Lie action, including
the prescribed scalar on the center, gives the comparison.  Each
exponential terminates on a nilpotent test ideal.  Changing the loop
lift by $K$ has precisely the Harish--Chandra equivariance already
imposed.  First-order differentiation yields the second line of
\eqref{eq:L1}, while products of formal elements give all higher
compatibilities.  By Proposition~\ref{loop:crystal} these are
calculations in crystal categories.

Here plots need not be smooth over a Schubert bound.  We record why
! plots still suffice.  For a finite-dimensional scheme $V$, a proper
map $\widetilde V\to V$ that is an isomorphism off a closed subset
$D$ gives a descent square with $D$ and
$\widetilde D=\widetilde V\times_VD$.  The open/closed triangles
prove this: open objects extended by star and closed objects have
the indicated pullbacks, and their extension maps agree by proper
direct-image/!-pull adjunction and Kashiwara's theorem.  These
identities remain valid with coefficient categories.  Resolution
of singularities and induction on the dimension of the closed
complement therefore reduce descent to smooth plots.  To compare
with any other smooth plot, resolve the closure of its lift over
the open locus and apply the same square; plots supported entirely
over the closed complement are handled by the induction.  This
also proves descent for morphisms and coherent data.  Twists and
restriction to affine plots do not change the argument.

The kernel formula is compatible with changes of level by the
projection formula.  For multiplication choose an intermediate
depth so that its left action on the second bounded right
quotient is trivial.  Project both inputs to these finite
quotients.  Pulling $M^g$ along the multiplication map gives the
iterated formula: on product plots conjugation composes, the
Maurer--Cartan term obeys its product rule, and the determinant
line multiplies.  Relative base change now identifies convolution
followed by the action with the two successive actions.
Coefficient functors commute with this integration because they
are continuous and parameter-linear.  The calculation with
strings proves coherence, and the delta kernel acts identically.

It remains to identify strong equivariance, including its derived
structure.  Take a deeper open $K_1$ normal in $J$, on which
$\chi_J$ vanishes.  The pro-unipotent forgetful identifications
reduce the calculation to the finite-dimensional smooth quotient
$Q=J/K_1$.  On a fixed input level, averaging and the bar maps
factor through such quotients by the kernel formula; further
output levels give the same calculation.

First replace $\ell$ by $j$ and work at a finite group stage of
$J$.  Harish--Chandra modules for $(j,K_1)$ are equivalent to
weakly right $J$-equivariant relative D-modules on $K_1\backslash J$.
Here weak equivariance is algebraic equivariance, without
identifying its differential with the D-module differential.
Indeed, translation descent of the underlying quasi-coherent
module identifies it with a representation of the stabilizer
$K_1$.  Differential-operator induction adds the monad
\[
                  U(j)\otimes_{U(\operatorname{Lie}K_1)}(-).
\]
Invariant first-order differentiation gives the Lie operators;
the stabilizer derivatives are already prescribed.  Brackets
and relative PBW identify the entire monad.  This is a derived
identification, using equivariant descent and its flat filtered
induction monad.

In this model the $Q$-action is strong translation, and its
conjugation kernel is exactly \eqref{eq:L1}.  Strong translation
invariants descend the D-module to the base while retaining the
weak $J$-action.  They therefore give algebraic representations
of $J$, as required.  A character twists this same computation.
Passing to pro $J$ is legitimate on induced generators: only a
finite-dimensional Lie quotient is added, and a generator
trivial on an open normal stage remains at that stage under
induction.  The representation cobar calculation proves the
colimit assertion, and taking $Q$-invariants commutes with it
by Proposition~\ref{loop:relative-equivariance}.

Finally add the remaining Lie operators from $\ell$ by their
induction monad.  Induction and its continuous forgetful right
adjoint are $Q$-linear.  For restriction this is immediate from
\eqref{eq:L1}; for induction the adjoint change on the
enveloping factor gives the same derivative rule.  Forgetful
is conservative by testing the induced generators.  Equivariant
bars thus give the Harish--Chandra induction monad for $(\ell,J)$.

There is no assumption here that infinite-level equivariance has
already been identified.  For a pro-unipotent congruence $K$, its
convolution idempotent is the identity on a level already
$K$-equivariant.  On a smaller level pass to a normal smaller
one; the finite calculation shows that the image of the
idempotent is exactly the modules equivariant for $K$.  This
identifies its image in the colimit defining $\mathcal M$.
All other compact opens are then handled by the same finite
quotient calculation.  This proves the asserted identification
and its compatibility with the forgetful functors.
\end{proof}

The spherical category constructed from the inner $P_0$ presentation
is the ordinary Kazhdan--Lusztig category.  Trivializing $P_0$ on the
disks gives the identification; changes of trivialization descend
because the modules are strongly positive-loop equivariant and the
extension splits there.  The same change of frames identifies the
localization constructions.

\subsection{The spherical--Whittaker pairing and its adjoints}

Set
\[
 \begin{split}
 C_S&=\mathcal M^H,\qquad \mathbf1_S=D(S),\\
 a_S&=\mathcal A\otimes_{D_*(H)}\mathbf1_S
           =D_{\mathcal L^{-s}}(LG/H),\\
 d_S&=(a_S)_{LN,\chi},\qquad
 Q_S=\mathcal M_{LN,\chi}.
 \end{split}
\]
The Grassmannian category uses closed ind-proper bounds.  Since $LN$
is an increasing union of compact pro-unipotent groups, its
coinvariants are the colimit of compact-subgroup coinvariants.
After the norm identification at each compact stage, the transition
functors are star averaging with character.  Denote the image of a
Grassmannian object $K$ in $d_S$ by $\mathbf d(K)$.

These conventions give
\begin{equation}\label{loop:whittaker-coinvariants}
                     d_S=\mathcal W_{-s}^-(S).
\end{equation}
The sign can be
checked without a choice of local coordinates.  Transpose the
Grassmannian pairing, which is star integration of ! tensor.
Coinvariants at phase $\chi$ become ! equivariance at phase $+\chi$
on the opposite, exponent-$s$ Grassmannian.  The model comparison
\eqref{eq:W2} therefore identifies $\mathbf d(K)$ with the functional
\begin{equation}\label{eq:loop-whittaker-functional}
                 W\longmapsto
                  \Gamma_{\mathrm{dR}}(K\otimes^!\mu^!W).
\end{equation}
The pairing is relative over $S$ when parameters are present.

For clarity, this duality also permits coefficient categories.
Grassmannian duality is computed by closed bounds and ordinary
D-module duality.  Compact-subgroup equivariant categories are
retracts by pro-unipotent averaging.  To pass from their colimit to
the tensor-compatible dual category, one may use the compact opens
of Proposition~\ref{walg:adolescent}: its presentation and transition
calculation apply to $a_S$, since deep-open-equivariant categories
generate it.
For a chosen compact open, use a cofinal system of closed bounds
stable under that group.  Such a system exists: on a fixed bounded
subscheme a sufficiently deep normal congruence acts trivially,
so its saturation factors through a finite-type quotient of the
compact open and remains bounded; take its stable closed closure.
On these bounds the open-equivariant stage uses finite-type
quotients with unipotent stabilizers.  The closed-bound transitions,
and then the compact-open transitions in that calculation, are
fully faithful with continuous $D(S)$-linear right adjoints.
Their stages are dualizable and their colimit remains dualizable.
Consequently its dual is the limit of the !-equivariance diagrams,
also after tensoring with coefficients.  This is the Whittaker
category in \eqref{eq:W2}.  Alternatively, the same conclusion
follows from the linear parametrized Whittaker comparison.

The loop action on a spherical module factors through the
Grassmannian.  Passing to $LN$-coinvariants gives the pairing
\begin{equation}\label{eq:L3}
                     F_S:C_S\otimes_{D(S)}d_S\longrightarrow Q_S.
\end{equation}
It is balanced over the spherical monoidal category
$\mathcal S_S=\operatorname{End}_{\mathcal A}(a_S)$, with factor
orientations chosen for the left loop action.  All its categories
and maps have the disjoint-cluster product property.  For $d_S$
this property comes from coinvariants, or equivalently from the
Whittaker model comparison with its dual pairing; this specifies
the normalization even when stack dimensions contribute shifts.

\begin{proposition}\label{loop:pairing-adjoints}
At the irrational levels under consideration, the balanced functor
\[
               \overline F_S:
               C_S\otimes_{\mathcal S_S}d_S\longrightarrow Q_S
\]
is fully faithful and has a continuous $D(S)$-linear right adjoint.
The projection
$C_S\otimes_{D(S)}d_S\to C_S\otimes_{\mathcal S_S}d_S$
also has a continuous $D(S)$-linear right adjoint.
These assertions hold over the full parameter scheme, including
its collision loci.
\end{proposition}

\begin{proof}
First identify the spherical category.  On any nonunit Schubert
orbit the twisting character of a stabilizer is nonintegral on
a one-dimensional torus.  At a nonzero coweight it is the
Killing pairing with that coweight, up to the fixed sign, times
the irrational exponent.  The corresponding equivariant
D-module category vanishes.  On a bounded Schubert chart this
argument applies to each distinct-point stratum; ! restrictions
to the collision strata give the same calculation for their
merged coweights.  After removing the unit locus there is no
remaining stratum.  Thus
\[
                         \mathcal S_S\simeq D(B_SH),
\]
with ! tensor, the convolution product on the unit orbit.

The multiplication of $D(B_SH)$ is ! pull along its diagonal.
Compute it using a finite smooth quotient of $H$ with
pro-unipotent kernel.  The diagonal is then smooth,
representable, and affine; its right adjoint is star direct
image with the smooth shift.  The ! projection formula makes
this adjoint bimodule-linear and continuous.  This uses the diagonal,
not the map $B_SH\to S$, and makes no compactness assertion about
the unit of $D(B_SH)$.  The balanced
projection for two $\mathcal S_S$-modules is obtained by
tensoring this multiplication and its adjunction with those
modules.  Hence its right adjoint has the same properties.
It is also conservative: the terms coming from the unbalanced
tensor generate the balanced tensor, as its tensor bar shows.

Put
\[
        a^r=\mathbf1_S\otimes_{D_*(H)}\mathcal A,
        \qquad
        \mathcal S_S=a^r\otimes_{\mathcal A}a_S.
\]
Consider the evaluation
\[
             e:a_S\otimes_{\mathcal S_S}a^r\longrightarrow\mathcal A.
\]
Before balancing, its geometric calculation is the evaluation
kernel for $H$ from Proposition~\ref{loop:relative-equivariance},
followed by
\[
             \mathcal A\otimes_{D_*(H)}\mathcal A
                            \longrightarrow\mathcal A.
\]
The latter functor is multiplication on $LG\times^H LG$.
Separating out the product coordinate identifies its fibers
with the ind-proper affine Grassmannian.  The right adjoint is
! pull and is continuous and bimodule-linear by base change.
Precisely, at each congruence quotient of the product coordinate
use closed proper Grassmannian bounds and then their compatible
system.  This constructs the adjoint in the pro convention;
no finite-type adjunction for the whole loop group is being
assumed.

Consequently the pre-balance evaluation has a continuous
bimodule-linear right adjoint.  The balanced evaluation $e$
has those same properties.  To check this formal passage,
write $p$ for balancing.  The right adjoint of $ep$ is
$p^Re^R$.  Since $p^R$ is continuous, conservative, and
bimodule-linear, it detects both the colimit comparison maps
for $e^R$ and its canonical linearity maps.  They are
isomorphisms because they are so for $(ep)^R$.

The functor $e$ is fully faithful.  Put
$E=a_S\otimes_{\mathcal S_S}a^r$.  Proposition~\ref{loop:relative-equivariance}
expresses $H$-invariants as a relative tensor operation and gives
\[
 E^{H\times H}\simeq
    a_S^H\otimes_{\mathcal S_S}(a^r)^H
       \simeq\mathcal S_S\otimes_{\mathcal S_S}\mathcal S_S,
 \qquad \mathcal A^{H\times H}\simeq\mathcal S_S.
\]
On these two-sided invariants evaluation becomes
$\mathcal S_S\otimes_{\mathcal S_S}\mathcal S_S\to\mathcal S_S$,
which is the identity.  Both adjoints are bimodule-linear,
so the unit of the adjunction is an isomorphism on these
invariants.  They generate the domain under the two loop
actions.  Indeed, $a_S$ and $a^r$ are generated by the loop
actions on their equivariant Grassmannian delta units, from
their relative tensor presentations; the tensor bar then
gives the assertion for the balanced product.  Continuity
and bimodule-linearity extend the unit isomorphism to the
whole domain.

Tensor this fully faithful adjunction with $\mathcal M$ over
$\mathcal A$ and take $LN$-coinvariants on the other side.
Both operations preserve the adjunction and its unit
isomorphism because both functors in the adjunction are
continuous and linear.  The resulting functor is precisely
$\overline F_S$.  This proves the first assertion, and the
earlier balancing calculation proves the second.  Every
calculation took place over $S$ and in its finite/pro
presentations, so the argument applies across collisions.
\end{proof}

\begin{corollary}[Parameter operations for the pairing]
\label{loop:parameter-operations}
Let $f:T\to S$ be a change between the smooth affine parameter charts
used here, including a principal open or a smooth closed collision or
graph chart.  For each of the families
$\mathcal E=\mathcal M,C,a,d,Q$, there is an identification
\[
             \mathcal E_T\simeq
                   D(T)\otimes_{D(S)}\mathcal E_S.
\]
The pairing and its continuous linear right adjoint base change under
these identifications.  For a smooth closed inclusion $i:Z\to S$ and a
principal open inclusion $j:U\to S$, they commute with
$i_*,i^!,j_*,j^!$, including the units and counits of the adjunctions.
\end{corollary}

\begin{proof}
For the Harish--Chandra categories, first apply the underlying
quasi-coherent functor $b$.  The monic divisor algebras and lattice
presentations commute with the indicated base changes.
Proposition~\ref{loop:hc-renormalization} identifies the mappings
between the base-changed compact generators, which still generate.
All the maps in question have finite Tor dimension since the source
and target charts are smooth.  The underlying category of
$D(T)\otimes_{D(S)}\mathcal E_S$ is
\[
 \operatorname{QCoh}(T)\otimes_{D(S)}\mathcal E_S
 \simeq\operatorname{QCoh}(T)\otimes_{\operatorname{QCoh}(S)}
                   \mathcal E_S^0,
\]
which is the category just computed.  The formal-diagonal transport
of Proposition~\ref{loop:crystal} restricts along $f$ and agrees on
every multiple infinitesimal diagonal with the transport on that
category.  This identifies the crystals.  The equivalence between
crystals and $D(S)$-module categories recalled at the start of this
section identifies their sections as asserted; see
\cite[Theorem~2.6.3 and Lemma~3.2.4]{Gaitsgory1Affine}.
For $a$ the assertion is the relative D-module tensor-product formula;
for $d$ and $Q$ it follows by base change in their coinvariant bars.

Tensoring both functors of a continuous $D(S)$-linear adjunction with
$D(T)$ tensors its unit and counit as well.  Thus its right adjoint
over $T$ is the base-changed right adjoint.  Finally, on any of these
module categories, the open and closed localization projectors
$j_*j^!$ and $i_*i^!$ are the actions of $j_*\omega_U$ and
$i_*\omega_Z$, respectively.  Linearity makes both adjoints commute with these
projectors.  Identifying their images with the restricted categories
gives the claimed commutation with the pushes and pulls, and the
tensor construction preserves their adjunction maps.
\end{proof}
